% !TeX root = ../navier-stokes-manuscript.tex
% ============================================================
% 2.2 Finite-Time Response Escalation
% ============================================================

\subsection{What finite time implies after one response escapes}\label{sub:FiniteTimeResponseEscalation}

The temporal Tower is already present once the same velocity field is differentiated in time.
Finite time adds a separate implication after one of those responses is known to escape: an integrable next response would give the escaping row a finite terminal value.
The argument below propagates an escape through existing rows under its stated hypotheses; it does not produce those rows.

Fix one Banach space \(B\), retain the Eulerian rows
\[
  V_n(t):=\partial_t^n u(\cdot,t),
\]
and suppose that \(V_n:[0,T_*)\to B\) is locally absolutely continuous, with Bochner derivative \(V_{n+1}\).
For \(0\le s<t<T<T_*\), the Banach-valued fundamental theorem of calculus gives
\[
  V_n(t)-V_n(s)
  =
  \int_s^t V_{n+1}(\tau)\,d\tau
\]
and
\[
  \lVert V_n(t)-V_n(s)\rVert_B
  \le
  \int_s^t \lVert V_{n+1}(\tau)\rVert_B\,d\tau.
\]
If
\[
  \sup_{0<t<T_*}\lVert V_n(t)\rVert_B
  =
  \infty,
\]
then the finiteness of \(\lVert V_n(0)\rVert_B\) forces
\[
  \int_0^{T_*}\lVert V_{n+1}(\tau)\rVert_B\,d\tau
  =
  \infty.
\]
Because \(T_*<\infty\), this in turn forces
\[
  \sup_{0<t<T_*}\lVert V_{n+1}(t)\rVert_B
  =
  \infty.
\]
This iteration is valid only while the same fixed space \(B\), absolute continuity, and the Bochner-derivative identity hold at every successive row.
That is the precise statement of finite-time response escalation for the temporal rows used in the paper.
